\section{Верификация и емпирична оценка}
\label{sec:verification_evaluation}

Тази глава обединява двата вида доказателства, върху които се опира настоящата работа. От една страна стоят ограниченията по време на компилация, модулните (модулни) тестове, интеграционните тестове и проверките на целевите системи на живо, чрез които се валидира коректността на архитектурата. От друга страна стои емпиричната оценка на асинхронния слой, чиято цел е да покаже кога базираното на корутини изпълнение носи реална практическа полза и кога неговата цена остава сравнима или по-висока от синхронната алтернатива.

Важният методологичен принцип е, че тези два вида доказателства не се разглеждат като взаимозаменяеми. Тестовата верификация отговаря на въпроса дали библиотеката е коректна спрямо собствените си договори (договорs) и ограничения на целевата система. Оценката чрез сравнителни тестове (сравнителен тест) отговаря на различен въпрос: носи ли асинхронният модел на изпълнение измерима полза при натоварвания, в които фазата на изчакване (фаза на изчакване) доминира над локалната изчислителна работа. Само съвместното присъствие на двата пласта прави архитектурните претенции на разработката достатъчно убедителни.

\subsection{Стратегия на верификацията}

Верификационната стратегия следва самата архитектура на библиотеката. Слоят по време на компилация се проверява чрез концепти, ограничения на възможностите и типизирано изграждане на междинно представяне на заявките. Слоят по време на изпълнение се проверява чрез модулни тестове за рендиране, параметризиране, нишкова безопасност, миграции, корутинна инфраструктура и асинхронен достъп. Най-външният слой се валидира чрез интеграционни тестове с реални целеви системи, така че да не се смесват коректност на имитациите (коректност на имитациите) и реално поведение на драйверите.

В текущата ревизия на хранилището пълният автоматизиран тестов набор включва 271 теста при изпълнение чрез \code{ctest}. Това число само по себе си не е научен резултат, но е важно като индикатор, че изложените в предходните глави механизми не са оставени на ниво концептуална схема. Те са покрити от систематични проверки по слоеве и по семейства целеви системи.

\begin{table}[H]
\centering
\small
\caption{Основни източници на верификационни доказателства}
\label{tab:evaluation_evidence_matrix}
\begin{tabularx}{\textwidth}{L{3.1cm}L{5.0cm}Z}
\toprule
\textbf{Източник} & \textbf{Какво валидира} & \textbf{Представителни артефакти} \\
\midrule
Ограничения по време на компилация и междинно представяне & типова коректност на заявки, базирано на възможности ограничаване и договори на конекторите & \path{tests/модулни/test_заявка.cpp}, \path{tests/модулни/test_postgresql_конектор.cpp}, \path{tests/модулни/test_mongodb_конектор.cpp} \\
Асинхронна инфраструктура & \code{Task<T>}, \code{thread\_pool}, \code{IoContext}, \code{async\_db<DB>}, съобразено с корутини поддържане на пул и транзакции & \path{tests/модулни/test_task.cpp}, \path{tests/модулни/test_нишка_pool.cpp}, \path{tests/модулни/test_io_context.cpp}, \path{tests/модулни/test_асинхронен_конектор.cpp} \\
Миграции и нишкова безопасност & логика за разлики в схемата (schema diff), семантика на отмяната (отмяна semantics), предпазители на връзки и конкурентна безопасност & \path{tests/модулни/test_schema_migration.cpp}, \path{tests/модулни/test_нишка_safety.cpp} \\
Интеграции с целеви системи на живо & поведение върху реални SQLite, PostgreSQL, MySQL, MongoDB, Redis, Neo4j и Cassandra среди & \path{tests/интеграционни/test_sqlite.cpp}, \path{tests/интеграционни/test_postgresql_live.cpp}, \path{tests/интеграционни/test_mongodb_live.cpp}, \path{tests/интеграционни/test_redis_live.cpp} \\
Емпирична оценка & пропускателна способност при симулирана фаза на изчакване и изолирани измервания на режийните разходи на корутините & \path{сравнителен тестs/bench_асинхронен.cpp} \\
\bottomrule
\end{tabularx}
\end{table}

Таблица~\ref{tab:evaluation_evidence_matrix} показва, че доказателствената база е хетерогенна по необходимост. Една библиотека за обектно-релационно съпоставяне по време на компилация не може да бъде защитена само с числа за пропускателна способност, както и асинхронна архитектура не може да бъде защитена само с проверки на концепти. Изискват се и двата вида доказателства.

\subsection{Каталог на модулните и интеграционните доказателства}

За да не остане стратегията за верификация само на равнище обща рамка, тя трябва да бъде разгърната и като конкретен каталог от артефакти в хранилището. Това е особено важно за настоящата работа, защото архитектурните претенции са разпределени между слоя на същностите, междинното представяне на заявките, договора на конектора, асинхронната инфраструктура и проверките на целевите системи на живо. Следващите таблици правят именно това: показват къде в хранилището са разположени основните доказателства и каква роля играе всяка тестова група.

\begin{table}[H]
\centering
\small
\caption{Основни модулни тестове за общите, основните и асинхронните подсистеми}
\label{tab:verification_unit_core}
\begin{tabularx}{\textwidth}{L{3.2cm}L{3.8cm}Y}
\toprule
\textbf{Файл} & \textbf{Подсистема} & \textbf{Какво валидира} \\
\midrule
\path{test_types.cpp} & базови псевдоними на типове & коректност на публичните псевдоними и базови помощни зависимости \\
\path{test_property.cpp} & слой за свойства на същностите & модел по време на компилация на скаларните полета и поведение на характеристиките на типовете \\
\path{test_relationship.cpp} & слой за връзки на същностите & логика за извеждане (извеждане) на връзките и поведение на стратегията за съхранение \\
\path{test_заявка.cpp} & междинно представяне на заявкитете и изразителен програмен интерфейс (изразителен API) & коректност на \code{select}, \code{insert}, \code{update}, \code{deleteq} и производните структури на заявки \\
\path{test_schema_migration.cpp} & миграционен слой & генериране на DDL, логика за разлики и поведение на автомат с крайни състояния \\
\path{test_нишка_safety.cpp} & помощници за конкурентност & поведение на пула от връзки, достъп, локален за нишката, и механизма на предпазителя на транзакцията \\
\path{test_task.cpp} & абстракция за корутинна задача & жизнен цикъл на \code{Task<T>} и базова семантика на изчакване/корутини \\
\path{test_нишка_pool.cpp} & планиране на работни нишки & публикуване, планиране и координация на ресурсите на работните нишки \\
\path{test_io_context.cpp} & асинхронна оркестрация & цикъл за събития и свързване на продълженията за асинхронния слой за изпълнение \\
\path{test_асинхронен_конектор.cpp} & фасада на асинхронния конектор & свързваща (bridge) логика между \code{async\_db<DB>} и куки на конектора, съобразени с корутини \\
\path{test_wire_protocol.cpp} & експериментален слой за изпълнение & ниско ниво на механизми по време на компилация / ориентирани към мрежовия протокол \\
\path{test_cpp26_отражение.cpp} & бъдеща езикова интеграция & ранен път за автоматизация, базирана на отражение (отражение) \\
\bottomrule
\end{tabularx}
\end{table}

Таблица~\ref{tab:verification_unit_core} е важна, защото показва, че модулният слой не се изчерпва с проверки, подобни на синтактичен анализатор на заявки. Той покрива едновременно логическия модел, механиката на изпълнение, миграциите и инфраструктурата, съобразена с корутини, върху която стъпва асинхронната архитектура.

\begin{table}[H]
\centering
\small
\caption{Модулни тестове, ориентирани към целевите системи, и техният фокус върху договора}
\label{tab:verification_unit_connectors}
\begin{tabularx}{\textwidth}{L{3.4cm}L{3.1cm}Z}
\toprule
\textbf{Файл} & \textbf{Целева система} & \textbf{Тип доказателство} \\
\midrule
\path{test_postgresql_конектор.cpp} & \code{PostgreSQLDB} & съответствие (compliance) с \code{is\_connector}, утвърждавания на възможностите и рендиране на SQL параметри \\
\path{test_mysql_конектор.cpp} & \code{MySQLDB} & договор за съединявания/транзакции, диалектно специфично рендиране и очаквания за свързване (свързване) \\
\path{test_mongodb_конектор.cpp} & \code{MongoDB} & негативно валидиране на възможности и рендиране на филтри, подобни на BSON \\
\path{test_redis_конектор.cpp} & \code{RedisDB} & отсъствие на съединявания/транзакции/агрегация и материализация на команди ключ-стойност \\
\path{test_cassandra_конектор.cpp} & \code{CassandraDB} & ограничения на възможностите и корелати за рендиране на CQL \\
\path{test_neo4j_конектор.cpp} & \code{Neo4jDB} & възможност за транзакции и материализация на заявки, ориентирани към графи \\
\bottomrule
\end{tabularx}
\end{table}

Тези модулни тестове, ориентирани към целевите системи, са критични, защото доказват не само позитивни възможности, но и честно декларирани ограничения. При библиотека с множество свързващи компоненти отрицателната верификация --- че дадена система \emph{не} поддържа дадена операция --- е също толкова важна, колкото и позитивната.

\begin{table}[H]
\centering
\small
\caption{Каталог на интеграционните тестове и тестовете на живо}
\label{tab:verification_integration_tests}
\begin{tabularx}{\textwidth}{L{3.2cm}L{3.0cm}Y}
\toprule
\textbf{Файл} & \textbf{Режим на изпълнение} & \textbf{Какво демонстрира} \\
\midrule
\path{test_mockdb.cpp} & винаги налична интеграционна базова линия & композиция на заявки, свързване на параметри, подготвени заявки и CRUD сценарии без външна целева система \\
\path{test_sqlite.cpp} & локален интеграционен тест & реален релационен път на изпълнение с утвърждавания на възможностите, съединявания, транзакции и попълване на резултати \\
\path{test_postgresql_live.cpp} & условен тест на живо & реална връзка, жизнен цикъл на таблици и изпълнение на CRUD върху PostgreSQL \\
\path{test_mysql_live.cpp} & условен тест на живо & реален път на MySQL драйвера, вмъкване/обновяване/изтриване и филтриране при избор \\
\path{test_mongodb_live.cpp} & условен тест на живо & изпълнение на сценарии, ориентирани към документи, срещу жива целева система \\
\path{test_redis_live.cpp} & условен тест на живо & материализация на ключ-стойност и хеш команди върху Redis среда \\
\path{test_cassandra_live.cpp} & условен тест на живо & изпълнение на широки колони и платформено-специфични CQL сценарии \\
\path{test_neo4j_live.cpp} & условен тест на живо & изпълнение на графови заявки през контролирана Docker среда \\
\bottomrule
\end{tabularx}
\end{table}

Таблица~\ref{tab:verification_integration_tests} показва, че доказателствата, базирани на хранилището, са разслоени. Не всички целеви системи имат еднаква експлоатационна цена за верификация на живо, но хранилището съдържа ясна стратегия как и кога подобни тестове се активират. Това е особено важно за тезата, защото избягва смесването между коректност на имитациите и реално поведение, подкрепено от драйвер.

\subsection{Карта на зрелостта според доказателствената опора}

Самият профил на възможностите не е достатъчен, ако липсват доказателства за реално поведение. Затова е полезно целевите системи да се разгледат и като стълбица на зрелост, определена не само от поддържаните операции, но и от наличната тестова опора. Тук зрелостта означава комбинация от коректност на договора, доказателства за изпълнение и експлоатационен реализъм.

\begin{table}[H]
\centering
\small
\caption{Свързващ компонент maturity карта според наличната доказателствена опора}
\label{tab:verification_backend_maturity_map}
\begin{tabularx}{\textwidth}{L{2.4cm}L{2.2cm}L{2.9cm}Y}
\toprule
\textbf{Свързващ компонент} & \textbf{Ниво на зрелост} & \textbf{Основна доказателствена опора} & \textbf{Интерпретация} \\
\midrule
\code{MockDB} & ниво 3 & \path{test_mockdb.cpp} & свързващ компонент-independent изпълнение базова линия за заявка договор-а и prepared заявка механиката \\
SQLite & ниво 3 към 4 & \path{test_sqlite.cpp} & локален реален изпълнение path с богата заявка evidence и възможност asserts \\
PostgreSQL & ниво 3 & \path{test_postgresql_конектор.cpp}, \path{test_postgresql_live.cpp} & силен договор плюс условна live свързващ компонент проверка \\
MySQL & ниво 3 & \path{test_mysql_конектор.cpp}, \path{test_mysql_live.cpp} & реален експлоатационни path с релационен профил и dialect-specific свързванеing \\
MongoDB & ниво 3 & \path{test_mongodb_конектор.cpp}, \path{test_mongodb_live.cpp} & document-oriented изпълнение evidence без SQL-style възможности \\
Redis & ниво 3 & \path{test_redis_конектор.cpp}, \path{test_redis_live.cpp} & key-value изпълнение evidence с ограничен, но честен договор профил \\
Neo4j & ниво 3 & \path{test_neo4j_конектор.cpp}, \path{test_neo4j_live.cpp} & graph изпълнение path с Docker-gated live verification \\
Cassandra & ниво 3 & \path{test_cassandra_конектор.cpp}, \path{test_cassandra_live.cpp} & широки колонни path с фокус върху формална пригодност и live coverage \\
\bottomrule
\end{tabularx}
\end{table}

Картата в Таблица~\ref{tab:verification_backend_maturity_map} е важна, защото предотвратява прекалено бинарно четене на свързващ компонент статуса. Конекторът не е просто наличен или липсващ; той може да бъде коректен по договор, валидиран при изпълнение, частично верифициран на живо или съобразен със схемата богат. Именно този по-нюансиран прочит прави верификационната глава по-научно убедителна и по-близка до реалния статус на хранилището.

\subsection{Методология на сравнителен тест оценката}

Емпиричната оценка е реализирана чрез Google Сравнителен тест~\cite{google_benchmark_user_guide} в \path{сравнителен тестs/bench_асинхронен.cpp}. Тя не цели да симулира пълния жизнен цикъл на конкретна база данни, а да изолира архитектурния въпрос дали базираното на корутини изпълнение освобождава работните нишки по-ефективно от синхронното блокиране, когато между две кратки CPU фази има интервал на изчакване.

Сравнителен тест наборът е разделен на две групи. Първата група измерва пропускателна способност на цели партиди от задачи. Синхронният базова линия изпълнява CPU работа, след което работна нишката блокира за предварително зададен интервал. Асинхронният вариант използва \code{SimulatedAsyncIO} изчакваем обект, който спира корутината, планира отложено възобновяване чрез опашка от таймери и връща продължаването обратно към \code{thread\_pool}. Втората група измерва изолирани режийни разходи: създаване и изчакване на \code{Task<int>}, \code{thread\_pool::post()} двупосочен път и вложени \code{co\_await} вериги.

Методологично е важно, че сравнителните тестове за пропускателна способност използват реално измерване, а не чисто CPU време. Причината е изрично отбелязана и в самия код на сравнителния тест: при корутинни сценарии с фаза на изчакване калибрирането по CPU време на Mac OS може да даде подвеждащо поведение, защото спряната задача почти не консумира CPU. Затова съпоставката се извършва чрез \code{UseRealTime()} и фиксиран брой итерации.

В настоящата глава се използват режим на издание резултатите от локалния \path{build-rel} дърво за изграждане. Те не трябва да се интерпретират като абсолютни производствени числа за всички възможни машини и драйвери. Тяхната роля е сравнителна: всички варианти се изпълняват върху една и съща машина, с една и съща параметрична мрежа, така че да се изведе относителният ефект от различния изпълнение модел.

\begin{table}[H]
\centering
\small
\caption{Семейства сравнителни тестове и тяхната цел}
\label{tab:benchmark_families}
\begin{tabularx}{\textwidth}{L{3.4cm}L{5.2cm}Y}
\toprule
\textbf{Семейство сравнителни тестове} & \textbf{Какво измерва} & \textbf{Основни параметри} \\
\midrule
\code{BM\_SyncThroughput} & пропускателна способност на блокираща базова линия, при която работната нишка остава заета по време на фазата на изчакване & брой нишки, брой задачи, латентност на изчакване, CPU итерации \\
\code{BM\_AsyncThroughput} & пропускателна способност на корутинен вариант със спиране/възобновяване и отложено възобновяване чрез опашка от таймери & брой нишки, брой задачи, латентност на изчакване, CPU итерации \\
\code{BM\_AsyncYieldThroughput} & цена на редуване на корутини (корутина interleave) без изкуствена латентност на изчакване & брой нишки, брой задачи, CPU итерации \\
\code{BM\_TaskCreateAndWait} & минимална цена на създаване и завършване на тривиална задача & единична корутина без външно планиране \\
\code{BM\_ThreadPoolPost} & цена на \code{post()} и двупосочен път (двупосочен път) на обещание/бъдеща стойност към работна нишка & размер на пула от нишки \\
\code{BM\_NestedCoAwait} & режийни разходи на вложено \code{co\_await} продължаване при различна дълбочина & дълбочина на корутинната верига \\
\bottomrule
\end{tabularx}
\end{table}

\subsection{Основни резултати за пропускателна способност}

Таблица~\ref{tab:throughput_results_release_mode} обобщава представителните резултати в режим на издание за най-информативните конфигурации. Показани са двойки, за които има директно сравнение между синхронния и асинхронния вариант при едни и същи параметри.

\begin{table}[H]
\centering
\small
\caption{Представителни резултати за пропускателната способност в режим на издание (режим на издание) за синхронно и асинхронно изпълнение}
\label{tab:throughput_results_release_mode}
\begin{tabularx}{\textwidth}{L{1.3cm}L{1.5cm}L{1.9cm}L{2.4cm}L{2.4cm}L{1.7cm}}
\toprule
\textbf{Нишки} & \textbf{Задачи} & \textbf{Латентност на изчакване} & \textbf{Синхронна пропускателна способност} & \textbf{Асинхронна пропускателна способност} & \textbf{Ускорение} \\
\midrule
4 & 100 & 10~$\mu$s & 194.6k ops/s & 658.5k ops/s & 3.38$\times$ \\
4 & 100 & 100~$\mu$s & 29.8k ops/s & 458.2k ops/s & 15.35$\times$ \\
4 & 100 & 1000~$\mu$s & 3.22k ops/s & 75.28k ops/s & 23.36$\times$ \\
4 & 1000 & 100~$\mu$s & 30.08k ops/s & 670.5k ops/s & 22.29$\times$ \\
8 & 100 & 100~$\mu$s & 57.13k ops/s & 258.5k ops/s & 4.53$\times$ \\
\bottomrule
\end{tabularx}
\end{table}

Най-важният извод е, че предимството на корутинния модел нараства не когато локалната CPU работа нарасне, а когато фазата на изчакване започне да доминира над нея. При 4 работни нишки и 100 задачи разликата между 10 и 1000 микросекунди ($\mu$s) латентност на изчакване увеличава относителното предимство на асинхронния вариант от 3.38$\times$ до 23.36$\times$. Това е именно очакваното поведение: при синхронната базова линия нишките се изчерпват от блокиране, докато при корутинния модел те се освобождават за друга работа.

Вторият важен извод е, че ефектът не изчезва при по-големи партиди. При 4 нишки, 1000 задачи и 100~$\mu$s латентност на изчакване асинхронният вариант достига 670.5k операции в секунда срещу 30.08k при синхронната базова линия. Това показва, че архитектурното предимство не е ограничено до „демо“ случай с малък брой задачи, а се запазва и когато конкуренцията за ресурси на работните нишки нарасне.

\begin{figure}[H]
\centering
\begin{tikzpicture}[x=1.95cm,y=0.24cm]
    \draw[->, thick] (0,0) -- (0,26) node[above] {\small ускорение асинхронен/синхронен};
    \draw[->, thick] (0,0) -- (6.4,0);

    \foreach \y in {5,10,15,20,25}
    {
        \draw[gray!35] (0,\y) -- (6.2,\y);
        \node[left] at (0,\y) {\scriptsize \y};
    }

    \foreach \x/\h/\labeltext/\valuetext in {
        0.7/3.38/{4 нишки\\100 задачи\\10~$\mu$s}/3.38,
        1.9/15.35/{4 нишки\\100 задачи\\100~$\mu$s}/15.35,
        3.1/23.36/{4 нишки\\100 задачи\\1000~$\mu$s}/23.36,
        4.3/22.29/{4 нишки\\1000 задачи\\100~$\mu$s}/22.29,
        5.5/4.53/{8 нишки\\100 задачи\\100~$\mu$s}/4.53}
    {
        \fill[green!28, draw=black] (\x-0.27,0) rectangle (\x+0.27,\h);
        \node[font=\scriptsize, align=center, above] at (\x,\h) {\valuetext$\times$};
        \node[font=\scriptsize, align=center, below=7pt] at (\x,0) {\labeltext};
    }
\end{tikzpicture}
\caption{Относително ускорение на асинхронната пропускателна способност спрямо синхронната базова линия при представителни конфигурации, ограничени от вход-изход (I/O-bound)}
\label{fig:async_speedup_cases}
\end{figure}

Фигура~\ref{fig:async_speedup_cases} визуализира същата тенденция в относителна форма. Дори когато броят на работните нишки се удвои до 8, асинхронният вариант остава приблизително 4.5$\times$ по-бърз при 100~$\mu$s латентност на изчакване. Това е важен практически приложим резултат: повече нишки намаляват част от натиска върху синхронната базова линия, но не елиминират структурното предимство на модела на изпълнение със спиране/възобновяване.

\subsection{Нулева латентност на изчакване (0ns) и режийни разходи за планиране на корутините}

При нулева изкуствена латентност на изчакване сравнението престава да измерва освобождаването на работните нишки от изчакване на вход-изход и започва да показва почти чистата цена на механиката за спиране/възобновяване. Поради това най-подходящата съпоставка, подкрепена от хранилището, комбинира \code{BM\_SyncThroughput} с нулева стойност за изчакване и \code{BM\_AsyncYieldThroughput}, който заменя базираното на таймер изчакване с еднократно прехвърляне (отстъпване) \code{PoolYield} обратно към \code{thread\_pool}.

\begin{table}[H]
\centering
\small
\caption{Представителни резултати в режим на издание (режим на издание) при нулева изкуствена латентност на изчакване}
\label{tab:zero_latency_yield_results}
\begin{tabularx}{\textwidth}{L{1.4cm}L{1.5cm}L{2.8cm}L{2.9cm}Y}
\toprule
\textbf{Нишки} & \textbf{Задачи} & \shortstack{\textbf{Синхронна пропуска-}\\\textbf{телна способност}} & \shortstack{\textbf{Пропускателна способност}\\\textbf{с отстъпване (асинхронен отстъпване)}} & \shortstack{\textbf{Отношение}\\\textbf{асинхронен/синхронен}} \\
\midrule
4 & 100 & 841.0k ops/s & 989.4k ops/s & 1.18$\times$ \\
4 & 1000 & 1.232M ops/s & 1.184M ops/s & 0.96$\times$ \\
8 & 100 & 810.9k ops/s & 514.4k ops/s & 0.63$\times$ \\
\bottomrule
\end{tabularx}
\end{table}

\begin{figure}[H]
\centering
\begin{tikzpicture}[x=1.65cm,y=4.0cm]
    \draw[->, thick] (0,0) -- (0,1.40) node[above] {\small пропускателна способност (M ops/s)};
    \draw[->, thick] (0,0) -- (5.8,0);

    \foreach \y/\label in {0.25/0.25,0.50/0.50,0.75/0.75,1.00/1.00,1.25/1.25}
    {
        \draw[gray!35] (0,\y) -- (5.6,\y);
        \node[left] at (0,\y) {\scriptsize \label};
    }

    \fill[blue!25, draw=black] (0.75,0) rectangle (1.05,0.841);
    \fill[green!28, draw=black] (1.15,0) rectangle (1.45,0.989);
    \node[font=\scriptsize, align=center, below=7pt] at (1.10,0) {4 нишки\\100 задачи};

    \fill[blue!25, draw=black] (2.55,0) rectangle (2.85,1.232);
    \fill[green!28, draw=black] (2.95,0) rectangle (3.25,1.184);
    \node[font=\scriptsize, align=center, below=7pt] at (2.90,0) {4 нишки\\1000 задачи};

    \fill[blue!25, draw=black] (4.35,0) rectangle (4.65,0.811);
    \fill[green!28, draw=black] (4.75,0) rectangle (5.05,0.514);
    \node[font=\scriptsize, align=center, below=7pt] at (4.70,0) {8 нишки\\100 задачи};

    \fill[blue!25, draw=black] (3.15,1.31) rectangle (3.33,1.37);
    \node[font=\scriptsize, anchor=west] at (3.37,1.34) {синхронно с 0 изчакване};
    \fill[green!28, draw=black] (4.35,1.31) rectangle (4.53,1.37);
    \node[font=\scriptsize, anchor=west] at (4.57,1.34) {асинхронно с отстъпване};
\end{tikzpicture}
\caption{Съпоставка на пропускателната способност при нулева изкуствена латентност на изчакване}
\label{fig:zero_latency_yield_throughput}
\end{figure}

Таблица~\ref{tab:zero_latency_yield_results} и Фигура~\ref{fig:zero_latency_yield_throughput} показват, че при липса на фаза на изчакване корутинният слой вече не получава структурния бонус, наблюдаван в режима, ограничен от вход-изход (I/O-bound). При 4 нишки и 100 задачи вариантът само с отстъпване (само с отстъпване) все още е приблизително 1.18$\times$ по-бърз, но при 4 нишки и 1000 задачи той е практически паритетен със синхронната базова линия (0.96$\times$), а при 8 нишки и 100 задачи пада до 0.63$\times$. Това е полезен коректив към силните резултати при ограничение от вход-изход: без външно изчакване корутинният механизъм не „създава“ пропускателна способност, а добавя малка, но измерима цена за планиране, която може да бъде компенсирана само частично.

\subsection{Интерпретация на микротестовете (microсравнителен тестs) за режийни разходи}

Изолираните измервания на режийните разходи (режийни разходи) показват, че корутинната инфраструктура сама по себе си не е доминиращият разход в сценариите, ограничени от вход-изход. При изграждане в режим на издание (издание) \code{BM\_TaskCreateAndWait} измерва приблизително 71.6~ns за създаване и завършване на тривиална задача. \code{BM\_ThreadPoolPost} остава около 5.0~$\mu$s независимо дали пулът има 1, 2, 4 или 8 работни нишки. \code{BM\_NestedCoAwait} достига около 324~ns при дълбочина 10 и около 2.375~$\mu$s дори при дълбочина 100.

Тези числа не трябва да се четат като пряка обещана цена на всяка реална заявка, но са важни с друго. Те показват, че режийните разходи на корутинния механизъм са порядъци по-малки от фазите на изчакване от 100--1000~$\mu$s, при които асинхронната пропускателна способност печели най-много. С други думи, наблюдаваното ускорение не е артефакт на измервателния инструмент, а следствие от това, че работните нишки престават да бъдат заключени в блокиращо изчакване.

\subsection{Ограничения на емпиричната оценка}

Емпиричната оценка съзнателно използва симулирана фаза на изчакване, а не пълно (от край до край) изпълнение срещу конкретен мрежов сървър за база данни. Това е ограничение, ако целта е да се предскаже абсолютна пропускателна способност за реална инсталация на PostgreSQL или Cassandra. Но то е и методологично предимство, ако целта е да се изолира моделът на изпълнение от влиянието на програмата за планиране на SQL (програма за планиране на SQL), мрежовите колебания, кеширането в драйвера и спецификите на ядрото за съхранение (ядро за съхранение).

Следователно настоящите резултати от сравнителните тестове не доказват, че „асинхронното е винаги по-бързо“. Те доказват по-тясна, но по-научно защитима теза: когато архитектурата работи в режим, ограничен от вход-изход (I/O-bound), и фазата на изчакване доминира над локалната изчислителна (CPU) работа, моделът, базиран на корутини, позволява съществено по-добро използване на ресурсите на работните нишки спрямо синхронната базова линия. Това е точно твърдението, което асинхронният слой на библиотеката трябва да обоснове.

В комбинация с тестовата верификация този резултат затваря двата основни въпроса на дипломната работа. От една страна библиотеката показва формална и тестово подкрепена дисциплина на коректност. От друга страна асинхронният модел на изпълнение показва измерима практически приложима полза именно в сценария, за който е проектиран. Така емпиричната оценка не стои встрани от архитектурата, а потвърждава нейния централен инженерен и научен смисъл.
